We ask whether self-supervised pretraining of a small CNN on an unlabeled pool of scikit-learn $8\times8$ digits improves a linear probe when only 10, 50 or 200 labels are available. We reimplement a SimCLR-style contrastive objective and rotation prediction, and compare them with supervised training from scratch of the same encoder, PCA plus logistic regression, raw pixels, and a frozen random encoder, using five seeds per budget on a fixed-size held-out test split. The SimCLR-style probe is the best system at every budget; its margin over supervised scratch is large at 10 labels, smaller at 50, and indistinguishable from zero at 200 labels (paired $p=\rhval{cmp/main/supervised-scratch/n200/test_acc/paired_p:2}$, five seeds). Rotation prediction fails in this configuration: it is below a random frozen encoder at 50 and 200 labels and not better than PCA at any budget (the differences from PCA are not resolved with five seeds); we hypothesise, but did not test, that the pretext task is solved without learning digit-discriminative features. Removing either augmentation family lowers SimCLR accuracy at 50 labels (photometric-only $p=\rhval{cmp/cmp_aug/simclr-photo-only/n50/test_acc/paired_p:3}$; geometric-only $-$\rhval{cmp/cmp_aug/simclr-geom-only/n50/test_acc/delta:3}, $p=\rhval{cmp/cmp_aug/simclr-geom-only/n50/test_acc/paired_p:2}$, not resolved). Accuracy is sensitive to temperature and pretraining length at about one standard deviation; $\tau=1.0$ was nominally higher (\rhval{sweep_tau/simclr-tau-1.0/n50/test_acc/mean:3} vs \rhval{main/simclr-style-probe/n50/test_acc/mean:3}) and accuracy saturates by 30 epochs. This is a small CPU study on one tiny dataset with reimplemented baselines and no hyper-parameter tuning; the conclusions should not be extrapolated to larger images or models.
